\section{Resultados y Conclusiones}

\subsection{Resultados}
La caja fue construida y ensamblada correctamente, cumpliendo los requerimientos de diseño y funcionalidad establecidos en la propuesta inicial. Se realizaron pruebas de objetos con diferentes etiquetas RFID. Vimos que la caja es capaz de detectar y registrar alrededor de unos 10 objetos antes de que el YRM se sature y deje de registrar de manera adecuada. Se realizaron pruebas de comunicación entre los microcontoladores y de la base de datos, en donde vimos que la información de los sensores y actuadores se almacena correctamente en la base de datos y se puede acceder desde nuestro dashboard web.

\subsection{Conclusiones}
La caja fue construida con el objetivo de poder almacenar y controlar el acceso al material de laboratorio, cumpliendo con los requerimientos básicos de diseño y funcionalidad mínimos para poder ser utilizada en un taller. Haciendo uso del dashboard, podemos corroborar que la información de que un objeto está en la caja o no, eficientando el control de inventario y acceso a materiales de una manera más eficaz; superando el 50\% de la capacidad máxima prevista para mejorar las operaciones. Ademas, con diez o menos objetos registrados, la caja funciona de manera adecuada, estando dentro del lumbral de saturación del YRM, estando adentro del 95\% de las etiquetas colocadas en la caja. Esto permite que el sistema sea escalable y pueda adaptarse a diferentes necesidades de almacenamiento y control de inventario.
